% SOURCE_AUTHORITY: sga5_fr_workpass.tex, lines 10465--10484 (LF-normalized unit).
% SOURCE_SHA256: 791F4EFFC5E02832D5D77ED03518C8156D6F07E4C8238B03545DB93D883FBB28
\emph{Demonstração.} O seguinte diagrama de implicações é evidente:
\[
\begin{array}{ccc}
\text{(ii)}&\Leftrightarrow&\text{(ii bis)}\\
\Downarrow&&\\
\text{(iii)}&\Leftrightarrow&\text{(iii bis)}
\end{array}
\qquad\qquad
\text{(i)}\Rightarrow\text{(i bis).}
\]
Demonstraremos separadamente que (i) $\Rightarrow$(ii bis) e (iii bis) $\Rightarrow$(i), o que será suficiente para concluir. A afirmação (i) $\Rightarrow$(ii bis) deduz-se sem dificuldade da sequência exata de Künneth
\[
0\to H^p(M)\otimes_V N\to H^p(M\otimes^L N)\to \operatorname{Tor}_1^V(H^{p+1}(M),N)\to0.
\]
A mesma sequência exata, para $N=k_V$, mostra que, se for cumprido (iii bis), tem
\[
H^p(M)\otimes k_V=0\quad(p\le m-1),\qquad
\operatorname{Tor}_1^V(H^m(M),k_V)=0.
\]
Daí se deduz (i) pelo lema de Nakayama e pelo critério local de planitude (EGA $0_{\mathrm{III}}$ 10.1.3).
